\chapter{Confiabilidade, Observabilidade e Opera\c{c}\~ao}
\label{chap:confiabilidade_observabilidade}

% ---------------------------------------------------------------
\section{Pr\'e-requisitos}
% ---------------------------------------------------------------

\begin{itemize}[leftmargin=*]
  \item Cap\'itulo~\ref{chap:pipeline_integrado_producao} --- pipeline
        completo em produ\c{c}\~ao com worker Lambda e DynamoDB;
  \item Familiaridade com JSON e logs estruturados.
\end{itemize}

% ---------------------------------------------------------------
\section{Escopo do cap\'itulo}
% ---------------------------------------------------------------

\begin{quote}\itshape
\textbf{Onde estamos na hist\'oria:} in\'icio do loop externo.
O sistema est\'a em produ\c{c}\~ao. Agora ele precisa ser observ\'avel:
toda resposta deve ter uma medida de confiabilidade, todo erro deve
ser rastreado, toda degrada\c{c}\~ao deve gerar alerta. Os cap\'itulos
11 a 14 instrumentam o loop externo que alimentar\'a o loop interno
com evid\^encias reais.
\end{quote}

Este cap\'itulo apresenta tr\^es camadas de observabilidade do sistema:

\begin{enumerate}[leftmargin=*]
  \item \textbf{\'Indice de confiabilidade:} m\'etrica estruturada
        calculada por resposta, baseada em cobertura de fontes RAG
        e sobreposi\c{c}\~ao l\'exica;
  \item \textbf{Rastreabilidade com OpenTelemetry:} traces e spans
        exportados ao Langfuse para depura\c{c}\~ao e auditoria;
  \item \textbf{Oper\-a\c{c}\~ao e resili\^encia:} retentativas,
        \textit{dead-letter queue} e monitoramento de estados no
        DynamoDB.
\end{enumerate}

No AWS Prescriptive Guidance, essa disciplina corresponde ao padr\~ao de
\textit{observer and monitoring agents}, com foco em telemetria, classifica\c{c}\~ao
de eventos e a\c{c}\~oes de escalonamento~\cite{awsagentic2025}.

% ---------------------------------------------------------------
\section{Objetivos de aprendizagem}
% ---------------------------------------------------------------

Ao concluir este cap\'itulo, o leitor ser\'a capaz de:

\begin{itemize}[leftmargin=*]
  \item Calcular e interpretar o \'indice de confiabilidade a partir
        das fontes RAG ativas;
  \item Configurar OpenTelemetry com exporta\c{c}\~ao ao Langfuse;
  \item Identificar as fontes de falha mais comuns e as estrat\'egias
        de resili\^encia do pipeline;
  \item Monitorar o ciclo de vida de uma requisi\c{c}\~ao usando os
        estados do DynamoDB.
\end{itemize}

% ---------------------------------------------------------------
\section{Conceitos te\'oricos}
% ---------------------------------------------------------------

\subsection{\'Indice de confiabilidade}
\label{sec:indice_confiabilidade}

Em sistemas RAG, uma resposta sem contexto suficiente \'e t\~ao
problem\'atica quanto uma resposta incorreta: o LLM pode alucinar
para preencher lacunas. O \'indice de
confiabilidade\index{\'indice de confiabilidade|textbf}
quantifica, por resposta, \textit{quanto do contexto necess\'ario
foi efetivamente recuperado}.

O c\'alculo \'e dividido em duas dimens\~oes:

\begin{itemize}[leftmargin=*]
  \item \textbf{Cobertura de fontes} (at\'e 0{,}95): cada camada RAG
        ativa contribui um peso espec\'ifico (Tabela~\ref{tab:pesos_fontes});
  \item \textbf{Sobreposi\c{c}\~ao l\'exica} (at\'e 0{,}15):
        fra\c{c}\~ao de palavras-chave do contexto que aparecem na
        resposta.
\end{itemize}

\begin{table}[htbp]
\centering
\caption{Pesos das fontes RAG no \'indice de confiabilidade.}
\label{tab:pesos_fontes}
\begin{tabular}{llr}
\toprule
\textbf{Fonte} & \textbf{Descri\c{c}\~ao} & \textbf{Peso} \\
\midrule
\texttt{opensearch\_bm25}    & BM25 sobre documentos e entrevistas & 0{,}40 \\
\texttt{neptune\_live}       & Grafo Neptune em tempo real          & 0{,}30 \\
\texttt{neptune\_sync}       & Grafo sincronizado (replica\c{c}\~ao) & 0{,}15 \\
\texttt{documentos\_knowledge} & \'Indice de pol\'iticas e FAQs     & 0{,}10 \\
\texttt{entrevistas\_corpus}  & Corpus de entrevistas de clientes   & 0{,}10 \\
\bottomrule
\end{tabular}
\end{table}

O n\'ivel resultante \'e classificado como: \texttt{alto}
($\geq 0{,}70$), \texttt{m\'edio} ($0{,}40$--$0{,}69$) ou
\texttt{baixo} ($< 0{,}40$).

\subsection{OpenTelemetry e Langfuse}

O \textit{OpenTelemetry} (OTel)\index{OpenTelemetry} \'e o padr\~ao
de observabilidade da CNCF para traces, m\'etricas e logs. O
\texttt{otel\_config.py} configura automaticamente:

\begin{itemize}[leftmargin=*]
  \item Exporta\c{c}\~ao de traces ao \textit{Langfuse Cloud} via
        OTLP/HTTP quando as vari\'aveis de ambiente
        \texttt{LANGFUSE\_PUBLIC\_KEY} e
        \texttt{LANGFUSE\_SECRET\_KEY} est\~ao presentes;
  \item Auto-instrumenta\c{c}\~ao das chamadas Anthropic via
        \texttt{openlit} (spans autom\'aticos para cada chamada ao
        Bedrock);
  \item \textbf{Modo silencioso:} se as chaves n\~ao estiverem
        configuradas, o m\'odulo importa sem erro ---
        \'util para desenvolvimento local.
\end{itemize}

% ---------------------------------------------------------------
\section{Implementa\c{c}\~ao orientada por passos}
% ---------------------------------------------------------------

\subsection{Passo 1 --- \'Indice de confiabilidade}

A fun\c{c}\~ao \texttt{\_confiabilidade} do \texttt{worker.py}
implementa o c\'alculo descrito na Se\c{c}\~ao~\ref{sec:indice_confiabilidade}:

\begin{lstlisting}[language=Python,
  caption={C\'alculo do \'indice de confiabilidade (worker.py).},
  label={lst:indice_confiabilidade}]
def _confiabilidade(bm25, sync, graph, resposta,
                    documentos="", entrevistas=""):
    fontes, cobertura = [], 0.0
    if bm25.strip():
        fontes.append("opensearch_bm25");  cobertura += 0.40
    if graph.strip():
        fontes.append("neptune_live");     cobertura += 0.30
    if sync.strip():
        fontes.append("neptune_sync");     cobertura += 0.15
    if documentos.strip():
        fontes.append("documentos_knowledge"); cobertura += 0.10
    if entrevistas.strip():
        fontes.append("entrevistas_corpus");   cobertura += 0.10

    ctx = " ".join(filter(None,[bm25,sync,graph,documentos,entrevistas]))
    overlap = 0.0
    if ctx and resposta:
        ctx_kw  = _palavras_chave(ctx)
        resp_kw = _palavras_chave(resposta)
        if ctx_kw and resp_kw:
            overlap = min(len(ctx_kw & resp_kw)
                          / min(len(ctx_kw), len(resp_kw)), 1.0)

    score = round(min(cobertura + overlap * 0.15, 1.0), 2)
    nivel = "alto" if score >= 0.70 else ("medio" if score >= 0.40
                                           else "baixo")
    return {"score": score, "nivel": nivel,
            "fontes_ativas": fontes,
            "detalhes": {"cobertura_fontes": round(cobertura,2),
                         "sobreposicao_lexica": round(overlap,2)}}
\end{lstlisting}

A resposta ao cliente inclui sempre o campo \texttt{indice\_confiabilidade},
permitindo que a camada de apresenta\c{c}\~ao decida se exibe um
aviso ao usu\'ario para respostas com n\'ivel \texttt{baixo}.

\subsection{Passo 2 --- Configurar observabilidade com OTel}

\begin{lstlisting}[language=Python,
  caption={Inicializa\c{c}\~ao do OTel no handler Lambda.},
  label={lst:otel_init}]
# Topo do worker.py, antes do handler
from otel_config import init as init_otel
init_otel()  # silencioso sem variaveis de ambiente
\end{lstlisting}

Para ativar o Langfuse, defina as vari\'aveis no Secrets Manager
e injete no ambiente Lambda via \texttt{terraform.tfvars}:

\begin{lstlisting}[language=bash,
  caption={Vari\'aveis de ambiente para Langfuse.},
  label={lst:langfuse_env}]
LANGFUSE_PUBLIC_KEY=pk-lf-...
LANGFUSE_SECRET_KEY=sk-lf-...
LANGFUSE_OTLP_ENDPOINT=https://cloud.langfuse.com/api/public/otel/v1/traces
\end{lstlisting}

\subsection{Passo 3 --- Estados da requisi\c{c}\~ao no DynamoDB}

O ciclo de vida de uma requisi\c{c}\~ao \'e rastreado por estados
no DynamoDB:

\begin{lstlisting}[language=bash,
  caption={Estados poss\'iveis de uma requisi\c{c}\~ao.},
  label={lst:estados_dynamodb}]
QUEUED      -> processamento iniciado
PROCESSING  -> agente em execucao
COMPLETED   -> resposta disponivel
FAILED      -> erro nao recuperavel
\end{lstlisting}

O endpoint \texttt{/status/\{request\_id\}} (\texttt{status.py})
consulta o DynamoDB e retorna o estado atual. Clientes implementam
\textit{polling} com intervalo exponencial at\'e \texttt{COMPLETED}
ou \texttt{FAILED}.

% ---------------------------------------------------------------
\section{Autoavalia\c{c}\~ao}
% ---------------------------------------------------------------

\begin{enumerate}[leftmargin=*]
  \item Uma resposta com \texttt{opensearch\_bm25} e
        \texttt{neptune\_live} ativos, sem sobreposi\c{c}\~ao l\'exica,
        tem \'indice de confiabilidade de quanto?\\
        \textbf{Gabarito:} $0{,}40 + 0{,}30 + 0{,}00 \times 0{,}15 = 0{,}70$
        $\to$ n\'ivel \texttt{alto}.

  \item O OTel exporta traces para o Langfuse mesmo quando
        \texttt{LANGFUSE\_PUBLIC\_KEY} n\~ao est\'a definido.\\
        \textbf{Gabarito: F} --- o \texttt{otel\_config.py} s\'o
        configura o exportador OTLP se ambas as chaves estiverem
        presentes.

  \item O estado \texttt{PROCESSING} no DynamoDB indica que a
        resposta j\'a est\'a dispon\'ivel para o cliente.\\
        \textbf{Gabarito: F} --- indica que o agente ainda est\'a em
        execu\c{c}\~ao; o cliente deve continuar o \textit{polling}.
\end{enumerate}

% ---------------------------------------------------------------
\section{Resumo}
% ---------------------------------------------------------------

\begin{itemize}[leftmargin=*]
  \item O \'indice de confiabilidade combina cobertura de fontes
        (at\'e 0{,}95) e sobreposi\c{c}\~ao l\'exica (at\'e 0{,}15).
  \item \texttt{opensearch\_bm25} e \texttt{neptune\_live} s\~ao as
        fontes de maior peso; sua aus\^encia sozinha j\'a classifica
        a resposta como n\'ivel \texttt{m\'edio}.
  \item OTel com Langfuse fornece visibilidade span-a-span das
        chamadas ao Bedrock.
  \item O modo silencioso do \texttt{otel\_config.py} garante que
        o c\'odigo funcione sem chaves Langfuse no ambiente de
        desenvolvimento.
  \item Estados DynamoDB permitem \textit{polling} ass\'incrono
        sem \textit{websocket}.
\end{itemize}

% ---------------------------------------------------------------
\section{Plano de testes do cap\'itulo}
% ---------------------------------------------------------------

\begin{itemize}[leftmargin=*]
  \item \textbf{TC-11-01:} chamar \texttt{\_confiabilidade} com apenas
        \texttt{bm25} ativo e verificar score $= 0{,}40$, n\'ivel
        \texttt{m\'edio}.
  \item \textbf{TC-11-02:} chamar com todas as 5 fontes ativas e
        sobreposi\c{c}\~ao 1{,}0 e verificar score $= 1{,}00$.
  \item \textbf{TC-11-03:} importar \texttt{otel\_config} sem
        vari\'aveis Langfuse e verificar que nenhuma exce\c{c}\~ao
        \'e lan\c{c}ada.
  \item \textbf{TC-11-04:} verificar que o endpoint
        \texttt{/status/\{id\}} retorna \texttt{404} para IDs
        inexistentes.
\end{itemize}

% ---------------------------------------------------------------
\section{Crit\'erios de aceite}
% ---------------------------------------------------------------

\begin{itemize}[leftmargin=*]
  \item \texttt{\_confiabilidade} passa todos os TC-11-0x.
  \item \texttt{otel\_config.init()} executa sem erro em ambiente
        local (validado por \texttt{validate\_sprint\_a.py}).
\end{itemize}

% ---------------------------------------------------------------
\section{Espa\c{c}o para expans\~ao iterativa}
% ---------------------------------------------------------------

\begin{itemize}[leftmargin=*]
  \item \textbf{v11.1:} dashboard CloudWatch com alarme para
        taxa de respostas n\'ivel \texttt{baixo} $> 10\%$.
  \item \textbf{v11.2:} integra\c{c}\~ao com PagerDuty para falhas
        em \textit{batch} cr\'itico.
  \item \textbf{Lacuna:} meta-judge autom\'atico --- usar o
        \'indice de confiabilidade como entrada do LLM Judge
        do cap\'itulo~\ref{chap:llm_as_judge}.
\end{itemize}
